\chapter{Support Vector Machine (Bonus-Vertiefung)}
\label{chap:svm}

\begin{myremark}[Algorithmusart in diesem Kapitel]
Funktionsart: \textbf{Klassifikation}. \\
Umsetzungsart: \textbf{selbstlernend}.
\end{myremark}

\begin{myremark}
Diese Bonus-Vertiefung ist \textbf{optional}. Sie setzt Kapitel~\cref{chap:einfuehrung} und Kapitel~\cref{chap:datenvorbereitung} voraus.
\end{myremark}

\section{Grundidee: Die optimale Trennebene}

\begin{mydefinition}[Support Vector Machine]
Eine \textbf{SVM} sucht eine Trennlinie (in 2D) bzw. Trennebene (in höheren Dimensionen),
die zwei Klassen mit möglichst \textbf{grossem Rand (Margin)} trennt. Dieser Parameter (Margin) wird allgemein durch die Regularisierungskonstante $C$ gesteuert.
Die Datenpunkte, die am nächsten an der Trennebene liegen, heissen \textbf{Support Vectors}.
\end{mydefinition}

\begin{figure}[H]
    \centering
    \foreach \margin in {0.1, 1.0, 10.0} {
            \begin{subfigure}[H]{0.32\textwidth}
                \includegraphics[width=\textwidth]{Grundlagen_Info/10_KI_und_Algorithmen/Figures/svm_boundary_\margin.pdf}
                \caption{Margin = \margin}
            \end{subfigure}
        }
    \caption{SVM-Entscheidungsgrenze mit Margin und Support Vectors für unterschiedliche Werte von $C$}
    \label{fig:svm-boundary-15}
\end{figure}

\section{Warum der Margin wichtig ist}

\begin{myexercise}[SVM intuitiv]
Betrachten Sie \cref{fig:svm-boundary-15}:
\begin{enumerate}
    \item Warum kann ein grosser Margin helfen, neue Datenpunkte besser zu klassifizieren?
          Erklären Sie in 3--4 Sätzen.
    \item Was passiert mit dem Margin, wenn ein Ausreisser direkt an der Grenze liegt?
\end{enumerate}
\end{myexercise}

\section{Soft Margin und der Parameter \texorpdfstring{$C$}{C}}

In der Praxis sind Daten selten perfekt trennbar.
Mit dem Parameter $C$ wird ein Kompromiss gesteuert:
\begin{itemize}
    \item \textbf{Grosses $C$:} enger Margin, wenige Fehler im Training (Overfitting-Risiko).
    \item \textbf{Kleines $C$:} breiter Margin, mehr Fehler erlaubt (Underfitting-Risiko).
\end{itemize}

\section{Kernel-Trick (konzeptuell)}

\begin{mydefinition}[Kernel-Trick]
Wenn die Klassen nicht linear trennbar sind, bildet der \textbf{Kernel-Trick}
die Daten in einen höher\-dimensionalen Raum ab, wo sie trennbar werden --
ohne die Transformation explizit zu berechnen.
Häufig verwendete Kernel: \textbf{RBF} (Radial Basis Function), Polynom.
\end{mydefinition}

\begin{mychallenge}[Kernel und Dimensionserhöhung]
Recherchieren Sie den RBF-Kernel und erklären Sie an einem Beispiel,
warum zwei Klassen, die im 2D-Raum nicht linear trennbar sind,
in einem höherdimensionalen Raum trennbar sein können.
\end{mychallenge}

\begin{mychallenge}[SVM vs.\ Random Forest]
Vergleichen Sie SVM und Random Forest bezüglich:
\begin{itemize}
    \item Interpretierbarkeit
    \item Skalierbarkeit (grosse Datensätze)
    \item Umgang mit vielen Features
\end{itemize}
Begründen Sie Ihre Aussagen mit Notebook-Experimenten.
\end{mychallenge}

\section{Gesellschaftliche Aspekte: Hyperparameter als verborgene Entscheidungen}

\begin{mydefinition}[Hyperparameter sind Wertsetzungen]
Der Kernel und die Regularisierungsparameter ($C$) sind nicht \enquote{optimal} im neutralen Sinne. Sie werden vom Menschen gewählt und spiegeln Wertsetzungen wider. Ein hoher $C$-Wert begünstigt Trainingsdaten-Accuracy, ein niedriger toleriert Fehler im Trainingsset -- es gibt keinen objektiven \enquote{besten\enquote} Wert.
\end{mydefinition}

\begin{myexample}[Kernel-Wahl und Fairness]
Bei einem Klassifikator für Hiring könnte die Kernel-Wahl entscheidend sein:
\begin{itemize}
    \item \textbf{Linearer Kernel:} Findet nur lineare Grenzen. Kann zu Underfitting führen, ist aber transparent (man kann die Gewichte interpretieren).
    \item \textbf{RBF-Kernel:} Sehr flexibel, kann komplexe Entscheidungsgrenzen finden. Aber diese sind völlig uninterpretierbar -- man kann nicht verstehen, warum ein Kandidat abgelehnt wurde.
\end{itemize}
\textbf{Frage:} Ist es fairer, ein transparentes Modell zu verwenden, das möglicherweise schlecht diskriminiert, oder ein genaues aber unerklärtes Modell? Die Antwort ist regulatorisch relevant (z.~B. GDPR Recht auf Erklärung).
\end{myexample}

\begin{myremark}[Notebook]
Praktische Umsetzung inkl. Übungen:
\href{\notebookbaseurl03_SVM_Classification.ipynb}{\texttt{03\_SVM.ipynb}}
\end{myremark}
